\lstset{style=Python}
\chapter{Anschluss von Datenbank an TigerJython-Game}
\section{Einführung}

In diesem Kapitel werden wir ein unser TigerJython-Game an eine Datenbank anschliessen, um neue Funktionalitäten wie ein Login-Fenster, einen Highscore und weitere Dinge zu ermöglichen. Hierzu möchten wir als erstes folgende Tabellen erstellen:

% Stil für alle Tabellen
\newtcolorbox{dbtable}[1]{colback=RoyalBlue!5!white,colframe=RoyalBlue!75!black,fonttitle=\bfseries,title=#1, width=7cm}

%\begin{figure}[H]
%\centering
%\begin{tikzpicture}
%    % Tabelle: Lieferant
%    \node[right] (scores) {
%        \begin{dbtable}{scores}
%            \begin{table}[H]
%            \centering
%            \begin{tabular}{cc|c}
%                        &Spalte&Typ\\\midrule\midrule
%                        \emoji{key}&\texttt{id} &\texttt{INTEGER}\\\midrule
%                        \emoji{up-right-arrow}  &          \texttt{user\_id}&\texttt{INTEGER}\node(score_userid){};\\\midrule
%                            &       \texttt{score}&\texttt{INTEGER}\\\midrule
%                              &      \texttt{date\_score}&\texttt{DATETIME}\\
%                        \end{tabular}
%            \end{table}
%        \end{dbtable}
%    };
%
%    % Tabelle: Kategorie
%    \node[right=3cm of scores] (users) {
%        \begin{dbtable}{users}
%            \begin{table}[H]
%                \centering
%                \begin{tabular}{cc|c}
%                    &Spalte&Typ\\\midrule\midrule
%                     \node(users_id){};\emoji{key}&\texttt{id} &\texttt{INTEGER}\\\midrule
%                      &          \texttt{user}&\texttt{VARCHAR(40)}\\\midrule
%                        &       \texttt{password}&\texttt{VARCHAR(40)}\\
%                    \end{tabular}
%                \end{table}
%        \end{dbtable}
%    };
%
%
%    % Beziehungen
%    \draw[->] (score_userid) -- (users_id);
%
%\end{tikzpicture}
%\end{figure}



\begin{figure}[H]
\centering
\begin{tikzpicture}[
    table/.style={
        matrix of nodes,
        nodes={anchor=west, minimum height=1em, text depth=0.2em, text height=1em, align=left},
        column 1/.style={right delimiter=|},
        row sep=0pt,
        column sep=2pt
    }
]

% Tabelle scores
\matrix (scores) [table] {
    & \textbf{Spalte} & \textbf{Typ} \\\midrule\midrule
    \emoji{key} & \texttt{id} & \texttt{INTEGER} \\\midrule
    \emoji{up-right-arrow} & \texttt{user\_id} & \texttt{INTEGER} \\\midrule
     & \texttt{score} & \texttt{INTEGER} \\\midrule
     & \texttt{date\_score} & \texttt{DATETIME} \\
};

% Tabelle users
\matrix (users) [table, right=5cm of scores] {
    & \textbf{Spalte} & \textbf{Typ} \\\midrule\midrule
    \emoji{key} & \texttt{id} & \texttt{INTEGER} \\\midrule
     & \texttt{user} & \texttt{VARCHAR(40)} \\\midrule
     & \texttt{password} & \texttt{VARCHAR(40)} \\
};

\node(userstitle)[above=5.5em of users.center, font=\bfseries, anchor=center]{Tabelle \texttt{users}};
\node(scorestitle)[above=7em of scores.center, font=\bfseries, anchor=center]{Tabelle \texttt{scores}};

% Pfeil von user_id (Zeile 3 in scores) zu id (Zeile 2 in users)
\draw [->,>=latex, thick] ([yshift=10pt]scores-3-3.east) -- +(40pt,0) -| ([xshift=-40pt,yshift=10pt]users-2-1.west) -- ([yshift=10pt]users-2-1.west);

\end{tikzpicture}
\end{figure}

Die Tabelle \texttt{users} speichert alle Benutzer-Daten für das Spiel, d.h., die Benutzer, welche sich in das Spiel eingeloggt haben. Die Tabelle \texttt{scores} speichert die Spielresultate jedes Spiels.

\begin{myexercise}[\texttt{users}-Tabelle erstellen]

Erstellen Sie ein Programm in TigerJython, um die Tabelle \texttt{users} in der Datenbank \texttt{game.db} zu erstellen. Die Datenbank wird mit dem unten stehenden Code automatisch angelegt, sollte sie nicht bereits existieren.
\begin{myattention}[Windows-Users]
\faWindows{} \textbf{Windows}: TigerJython im Admin-Modus ausführen, da sonst nicht auf die lokale SQLite-Datenbank zugegriffen werden kann!
\end{myattention}

\begin{lstlisting}
from sqlite3 import *
from prettytable import *   # Befehle zur Ausgabe von Tabellen

# Datenbank game.db (wenn nötig) erzeugen und verbinden
with connect("game.db") as con:
    # Datenbank-Cursor (=Lese/Schreibestift) holen
    cursor = con.cursor()
    
    #TODO Mehrzeiliger SQL-Code zur Erstellung einer Tabelle
    sql = """
        CREATE TABLE IF NOT EXISTS users (
           ...
    """
    
    # SQL ausführen
    cursor.execute(sql)

    print("New table X was created!")

\end{lstlisting}

\textbf{Tipp}: Primärschlüssel müssen erstellt werden mit dem Befehl \lstinline[style=sql]|PRIMARY KEY AUTOINCREMENT|:
\begin{lstlisting}[style=sql]
CREATE TABLE IF NOT EXISTS users(
	id INTEGER PRIMARY KEY AUTOINCREMENT, ...
)
\end{lstlisting}
Mit dem Zusatz \lstinline[style=sql]|AUTOINCREMENT| werden die Werte der Spalte \texttt{id} automatisch bei jedem Eintrag um 1 vergrössert, so dass kein Wert für \texttt{id} beim \lstinline[style=sql]|INSERT|-Befehl mitgegeben werden muss.
\begin{myanswer}
\begin{lstlisting}
from sqlite3 import *

# Verbindung zur DB game.db öffnen
with connect("game.db") as con:
    # Datenbank-Cursor (=Lese/Schreibestift) holen
    cursor = con.cursor()
    
    # SQL-Code zur Erstellung der Tabelle
    sql = """
        CREATE TABLE IF NOT EXISTS users (
            id_INTEGER PRIMARY KEY AUTOINCREMENT,
            user VARCHAR(40),
            password VARCHAR(40) 
        )
    """
    
    # SQL ausführen
    cursor.execute(sql)

    print("Table users was created!")
\end{lstlisting}
\end{myanswer}
\end{myexercise}



\begin{myexample}[SQL-Tabellen oder -Werte in Python ausgeben]
Resultate von SQL-Abfragen können auf zwei Arten ausgegeben werden:
\begin{enumerate}
\item Einzelne Werte ausgeben:
\begin{lstlisting}[style=Python]
cursor.execute("SELECT * FROM games")
resultset = cursor.fetchall()
# gib die dritte Spalte (2) der ersten Zeile (0) aus
print(resultset[0][2]) 
\end{lstlisting}

\item Um eine ganze Tabelle, schön dargestellt auszugeben, muss die Library \lstinline|prettytable| geladen und wie folgt verwendet werden:
\begin{lstlisting}[style=Python]
from prettytable import *

# ... 

cursor.execute("SELECT * FROM games")
printTable(cursor)
\end{lstlisting}
\end{enumerate}
\end{myexample}

\begin{myexercise}[Benutzer-Registrierung]
Schreiben Sie ein Programm, um Sie mittels \lstinline|input("...")| nach einem Benutzernamen und Passwort zu fragen. Beides sollte danach mit Insert in der Tabelle \texttt{users} abgespeichert werden. Überprüfen Sie am Ende das Resultat mit folgendem Code:
\begin{lstlisting}
sql = "SELECT * FROM users"
cursor.execute(sql)
printTable(cursor)
\end{lstlisting}
\textbf{Tipp}: Mit dem Text-Befehl \lstinline|.format()| lassen sich die Variablen-Werte innerhalb der Platzhalter \lstinline|{}| in einen Text einfügen.

\begin{lstlisting}
sql = """INSERT INTO users (name, password)
        VALUES('{}', '{}')""".format(username, password)
\end{lstlisting}
Die drei Anführungszeichen ermöglichen es, Texte über mehrere Zeilen zu schreiben.
\begin{myanswer}
\begin{lstlisting}
from sqlite3 import *
from prettytable import * # Befehle zur Ausgabe von Tabellen

# Datenbank game.db (wenn nötig) erzeugen und verbinden
with connect("game.db") as con:
    # Datenbank-Cursor (=Lese/Schreibestift) holen
    cursor = con.cursor()
    
    # Benutzername und Passwort abfragen
    username = input("Enter username")
    password = input("Enter password")
    
    # SQL-Code zum Einfügen eines neuen Datensatzes
    sql = """
        INSERT INTO users (user, password)
        VALUES
            ('{}', '{}')
    """.format(username, password)
    
    # SQL ausführen
    cursor.execute(sql)
    
    # Tabelle ausgeben
    sql = "SELECT * FROM users"    
    cursor.execute(sql)
    printTable(cursor)

\end{lstlisting}
\end{myanswer}
\end{myexercise}

\begin{mychallenge}[Sicherheits-Überprüfung]
Fügen Sie eine Bedingung ein, wonach eine Aufforderung erscheint, einen anderen Benutzernamen einzugeben, sollte es einen bestimmten Benutzer bereits geben.
\begin{myanswer}
\begin{lstlisting}
from sqlite3 import *

with connect("game.db") as con:
    # Datenbank-Cursor (=Lese/Schreibestift) holen
    cursor = con.cursor()

    # Benutzername abfragen
    repeat:
        username = input("Enter username")
        sql = """
            SELECT * 
            FROM users
            WHERE user = '{}'
        """.format(username)
        cursor.execute(sql)
        resultset = cursor.fetchall()
        if len(resultset) == 0:
            # Noch unbenutzer Benutzernamen --> alles OK
            break
        else:
            print("This username is already taken. Choose a different one!")
    
    password = input("Enter password")

    # Restlicher Code

\end{lstlisting}
\end{myanswer}
\end{mychallenge}

\begin{mychallenge}[Passwort verschlüsseln]
In der Praxis werden Passwörter nicht im Klartext in Datenbanken gespeichert, sondern mit einer kryptografischen Einwegfunktion, einem sogenannten \textit{hash}, verschlüsselt. Dies dient dazu, dass nicht einmal der Administrator einer Webseite, welcher Zugriff auf die Datenbank hat, die Passwörter der BenutzerInnen lesen kann. Sie können anstelle des Passworts die Ausgabe der SHA1-Hashfunktion des Passworts wie folgt speichern:
\begin{lstlisting}
import hashlib
password = 'helloworld!'              # Ihr Klartext-Passwort zum "hashen"
hash_object = hashlib.sha1()          # SHA-1-Hash-Objekt erstellen
hash_object.update(password.encode()) # Hash-Objekt updaten mit der Byte-Darstellung des Passworts
hex_dig = hash_object.hexdigest()     # Hexadezimale Darstellung erhalten
print("SHA-1 Hash:", hex_dig)         # Hash ausgeben

\end{lstlisting}
\begin{myanswer}
\begin{lstlisting}
from sqlite3 import *
import hashlib

# ...

    password = input("Enter password")    # Ihr Klartext-Passwort zum "hashen"
    hash_object = hashlib.sha1()          # SHA-1-Hash-Objekt erstellen
    hash_object.update(password.encode()) # Hash-Objekt updaten mit der Byte-Darstellung des Passworts
    password = hash_object.hexdigest()    # Hexadezimale Darstellung erhalten

    # Restlicher Code
\end{lstlisting}
\end{myanswer}
\end{mychallenge}

\begin{myexercise}[Users-Tabelle mit Game verknüpfen]
Erweitern Sie Ihr TigerJython-Game so, dass man zu Beginn aufgefordert wird, ein gültiges Login einzugeben. Erst wenn dies passiert ist, startet das Spiel. Sie können hierzu folgende Vorlage vervollständigen:
\begin{lstlisting}
while True:
    username = input("Enter username")

    # TODO Den User mit diesem Benutzernamen in der Tabelle users suchen
    
    if True: # TODO Bedingung anpassen: Wenn der User gefunden wurde…
        break
    else:
        print("Unknown user. Try again.")

while True:
    password = input("Enter password")

    if True: # TODO Bedingung anpassen: Wenn das eingegebene Passwort stimmt…
        break
    else:
        print("Incorrect password. Try again.")
\end{lstlisting}
\begin{myanswer}
\begin{lstlisting}
from sqlite3 import *

with connect("game.db") as con:
    cursor = con.cursor()
    
    repeat:
        username = input("Enter username")
    
        # Den User mit diesem Benutzernamen in der Tabelle users suchen
        sql = """
            SELECT * 
            FROM users
            WHERE user = '{}'
        """.format(username)
        cursor.execute(sql)
        resultset = cursor.fetchall()
        
        if len(resultset) > 0: # Wenn der User gefunden wurde...
            break
        else:
            print("Unknown user. Try again.")
    
    repeat:
        password = input("Enter password")
    
        if resultset[0][2] == password: # Wenn das eingegebene Passwort stimmt
            break
        else:
            print("Incorrect password. Try again.")
            
    print("Welcome", username)

\end{lstlisting}
\end{myanswer}
\end{myexercise}

\begin{myexercise}[\texttt{scores}-Tabelle mit Game verknüpfen]
Schreiben Sie ein Programm zur Erstellung der \texttt{scores}-Tabelle und erweitern Sie ihr Game so, dass am Ende jedes Spiels die Punktzahl darin abgespeichert wird.

\texttt{Tipp}: Verwenden Sie den SQL-Befehl \lstinline[style=sql]|DATETIME('NOW')|, um die aktuelle Datums- und Zeit-Angabe zu erhalten.
\begin{myanswer}
\begin{lstlisting}
# ...

with connect("game.db") as con:
    # Datenbank-Cursor (=Lese/Schreibestift) holen
    cursor = con.cursor()

   # Benutzername abfragen
    sql = """
        INSERT INTO scores (user_id, achieved_score, achieved_date)
        VALUES
            ({}, {}, DATETIME('NOW'))
    """.format(user_id, score)
 

\end{lstlisting}
\end{myanswer}
\end{myexercise}

\begin{myexercise}[Highscore-Ansicht hinzufügen]
Erweitern Sie Ihr Videogame mit einer Highscore-Ansicht, in welcher die zehn besten Punktestände inklusive Benutzernamen und Datum angezeigt werden.
\begin{myanswer}
\begin{lstlisting}
from sqlite3 import *

with connect("game.db") as con:
    cursor = con.cursor()
    
    #TODO TOP 10 scores inkl. Userdaten auslesen
    sql = """
        SELECT *
        FROM scores s
        JOIN users u ON u.id = s.user_id
        ORDER BY s.achieved_score DESC
        LIMIT 10
    """
    
    cursor.execute(sql)
    resultset = cursor.fetchall()
    
    for i in range(len(resultset)):
        hs = resultset[i]
        print(i+1, " User:", hs[5], "    Points:", hs[2])
        #TODO Mit GPanel-Befehl text() auf die Leinwand schreiben
\end{lstlisting}
\end{myanswer}

\end{myexercise}

\begin{mychallenge}
Erstellen Sie einen Super-User und verändern Sie ihr Spiel so, dass die Grafik verändert wird, wenn der Super-User eingeloggt ist (z.B. wird ein goldener Ball statt einem normalen Ball angezeigt). Lassen Sie Ihrer Kreativität freien Lauf!
\end{mychallenge}

\begin{mychallenge}
Zeigen Sie bei jedem Game-Start einen Text mit dem bisher höchsten Punktestand an, um den User dazu zu animieren, den Highscore zu knacken!
\end{mychallenge}